\chapter{Introducción}
\label{cap:introduccion}

% La introducción responde a ¿qué? y ¿por qué?: contexto, justificación, objetivos,
% alcances, método (a grandes rasgos) y estructura de la memoria.

\section{Contexto y motivación}
\label{sec:contexto}
% En las democracias representativas los partidos comunican sus valores mediante
% programas y discursos públicos, pero a menudo hay discrepancias entre lo que
% declaran y lo que hacen sus representantes. Por qué importa: confianza, transparencia.

\section{Descripción del problema}
\label{sec:problema}
% El desfase declarado/práctica está poco medido empíricamente; las fuentes están
% dispersas y en formatos heterogéneos; los estudios suelen mirar una sola dimensión.

\section{Objetivos}
\label{sec:objetivos}

\subsection{Objetivo general}
\label{sec:objetivo-general}
% Construir y analizar bases de datos públicas que permitan detectar y describir
% discrepancias entre el discurso programático de los partidos y la práctica
% parlamentaria de sus representantes.

\subsection{Objetivos específicos}
\label{sec:objetivos-especificos}
% 1. Recolectar y organizar fuentes (redes sociales, programas, datos parlamentarios).
% 2. Construir una base de programas partidarios comparable (MPDB).
% 3. Recolectar y estructurar información parlamentaria (intervenciones, votaciones).
% 4. Estandarizar y unificar las bases en un formato apto para análisis computacional.
% 5. Aplicar modelos de NLP (clasificación MARPOR, modelado de tópicos) sobre los corpus.
% 6. Identificar y cuantificar discrepancias entre lo declarado y lo revelado por el voto.

\section{Alcances}
\label{sec:alcances}
% Qué cubre y qué no: foco en Francia (XV legislatura 2017-2022); 5 corpus;
% análisis a nivel de partido; limitaciones de cobertura y de los datos.

\section{Estructura de la memoria}
\label{sec:estructura}
% Breve mapa de los capítulos siguientes.
